\section{\label{sec:Introduction} 引言}

随着 ATLAS 与 CMS 两个实验组在 13 TeV 运行中各自积累了超过 140 fb$^{-1}$ 的数据样本，大型强子对撞机（LHC）已进入精确物理时代。理论预言的改进也跟上了实验精度的步伐：目前已有若干对撞机过程的计算达到了 QCD 耦合强度展开的次次领头阶（NNLO）。这样的精度对于严格检验标准模型（SM）以及寻找超出标准模型（BSM）物理的迹象都是必需的，因为迄今为止尚未出现 BSM 物理的``确凿''信号。QCD 理论的精确预言需要相应精确的部分子分布函数（PDF）~\cite{Dulat:2015mca,Harland-Lang:2014zoa,Ball:2017nwa,Alekhin:2017kpj,Accardi:2016qay,Harland-Lang:2019pla,Bertone:2017bme,Manohar:2017eqh}，而这又要求在诠释 LHC 实验、提取关于 SM 乃至可能的 BSM 物理的重要信息方面不断进步。

为此，我们给出新一族 CTEQ-TEA 部分子分布函数，记作 CT18。这些 PDF 在 QCD 耦合常数 $\alpha_s$ 的次领头阶（NLO）和 NNLO 下给出。CT18 PDF 更新了文献~\cite{Dulat:2015mca} 中给出的 CT14 族 PDF。在新分析中，我们纳入了来自 ATLAS、CMS 与 LHCb 合作组、质心能量为 7 和 8 TeV 的多种新 LHC 数据，涵盖单举喷注、$W/Z$ 玻色子和正反顶夸克对的产生；同时，这一更新保留了此前 CT 全球 QCD 分析中关键的``传承''数据，例如 HERA I+II 合并深度非弹性散射（DIS）数据集，以及固定靶实验和费米实验室 Tevatron $p\bar p$ 对撞机上的测量。ATLAS 与 CMS 对相似运动学区域内过程的测量为数据提供了关键的交叉检验；LHCb 的测量则常常允许外推到其他实验未覆盖的新运动学区域。某些过程（如 $t\bar{t}$ 产生）允许多个可观测量的测量，它们为 PDF 的确定提供相似的信息。
除 PDF 本身之外，我们还给出相关的 PDF 光度，以及 LHC 标准烛光截面及其不确定度的预言。

CT18 分析的目标是尽可能覆盖每个测量的宽运动学范围，同时仍保持数据与理论的合理符合。例如对于 ATLAS 7 TeV 喷注数据~\cite{Aad:2014vwa}，若不引入 ATLAS 合作组提供的系统误差去关联~\cite{Aaboud:2017dvo}，就无法同时使用所有快度区间。即便做了该去关联，新喷注实验得到的 $\chi^2$ 仍不理想，导致对 PDF 的约束效力下降。ATLAS 和 CMS 都在两种不同的喷注半径 $R$ 取值下进行了单举截面测量；对两个实验我们都选择了较大 $R$ 值的数据，因为对此 NNLO（固定阶）预言应有更高精度。我们采用单举喷注横向动量作为 QCD 标度 $\Q\! =\! p_T^{jet}$ 来计算喷注截面预言，这与过去 NLO 下的用法一致。其结果与采用 $\Q\! =\! H_T$ 的类似计算基本一致~\cite{Currie:2016bfm,Currie:2017ctp,Currie:2018xkj}。

用于与全球拟合数据比较的理论预言均在 NNLO 完成：或通过 \texttt{fastNLO}~\cite{Britzger:2012bs,Wobisch:2011ij}、\texttt{ApplGrid}~\cite{Carli:2010rw} 等快速插值表结合 NNLO/NLO $K$ 因子间接实现，或（对顶夸克相关可观测量）通过 \texttt{fastNNLO} 网格~\cite{Czakon:2017dip,fastnnlo:grids} 直接实现。

在理想世界中，所有这些数据集应彼此完全相容，但实际观测到的差异确实导致数据集之间的若干张力以及相反方向的拉动。开展全球 PDF 分析的一个关键方面，就是处理那些给拟合带来张力的数据集，同时保住组合数据集对既有 PDF 约束的改进能力。在某些情况下，某数据集的张力可能大到必须将其从全球分析中剔除，或仅在一个独立迭代的新 PDF 组中纳入。

本文将描述我们发现需要特殊处理的高精度 ATLAS 7 TeV $W/Z$ 快度分布数据：如我们所见，它们倾向于增大低 $x$ 区的奇异夸克分布。特别地，尽管其他 PDF 分析组（如 MMHT，见文献~\cite{Thorne:2019mpt}）指出这些 ATLAS $W/Z$ 数据可以在与 CT18 相当水平的 $\chi^2/\mathrm{d.o.f.}$ 下拟合，我们发现这样的 $\chi^2$ 反映的是与我们全球分析中许多其他数据的系统性张力。此外，实验合作组常用的标准 Hessian 剖析（profiling）技术显著低估了把 ATLAS 7 TeV $W/Z$ 数据纳入 CT18 拟合后所能达到的最小 $\chi^2$。因此我们在一个替代拟合 CT18A 中单独处理这些测量，该拟合将在 Sec.~\ref{sec:alt} 中介绍。另一变体 CT18X 对低 $x$ DIS 截面的计算采用特殊标度 $\mu^2_{F,x} \equiv 0.8^2 \left(Q^2 + 0.3\mbox{ GeV}^2/x_B^{0.3}\right)$，该标度模拟低 $x$ 重求和的效应。这两种修改都会使低 $x$ 夸克与胶子分布增大。最后，这两个变体被合并为一个组合的替代拟合 CT18Z。由于 CT18Z PDF 与 CT18 差异最大，全文我们将展示大量基于 CT18 与 CT18Z PDF 的结果，而把与 CT18A 和 CT18X 的比较放到 Appendix~\ref{sec:AppendixCT18Z}，其中还提供了与 CT18Z 的更多深入比较。关于如何按用户需求在四个 PDF 组合中选择其一的建议见 Sec.~\ref{sec:Conclusions} 开头。

我们目前的全球分析分四个阶段进行。首先，使用 \texttt{PDFSense}~\cite{Wang:2018heo,Hobbs:2019gob}——一个基于 Hessian 方法~\cite{Pumplin:2001ct,Pumplin:2002vw} 对 QCD 数据进行快速普查的程序——来挑选预计对全球 PDF 组影响最大的数据集。这种挑选考虑了数据在给定 $x$ 区间内对特定 PDF 的敏感度，它既反映这些数据与给定 PDF 的相关性，也反映数据集规模及其统计与关联系统误差的大小。例如，对撞机单举喷注数据和顶夸克数据都与高 $x$ 胶子强相关，但单举喷注数据因数据点数目大得多而具有更大的敏感度。其次，用 \texttt{ePump}~\cite{Schmidt:2018hvu,Hou:2019gfw} 在 Hessian 近似下快速考察每个入选数据集的定量影响。第三，利用全部这些数据集完成完整的全球 PDF 拟合；近期对 CT 全球分析代码的增强大大提升了计算速度。最后，利用拉格朗日乘子（LM）方法~\cite{Stump:2001gu} 评估关键数据集对特定感兴趣运动学点上某些 PDF 以及 $\alpha_s(M_Z)$ 数值的影响。
为使参数化偏差最小，我们对 CT18 测试了不同参数化方案，例如对奇异夸克 PDF 采用更灵活的参数化。在某些运动学区域，数据对某些 PDF 的约束较少；特别地，我们施加了 LM 约束，把 $x < 10^{-5}$ 处的奇异性 PDF 限制在物理合理的取值，概括于 App.~\ref{sec:AppendixParam}。

本文结构如下。Sec.~\ref{sec:Datasets} 先给出 CT18 全球分析关键步骤与结果的概要总结，随后综述所选实验数据及 CT18 发布版中的替代拟合（CT18Z、CT18A 与 CT18X）。本节简明总结了大多数读者关心的关键结果；后续各节与附录则就具体方面和结论展开。

Sec.~\ref{sec:Theory} 详细介绍 CT 拟合方法学在理论/计算方面的更新，以及依赖具体过程的计算细节。Sec.~\ref{sec:OverviewCT18} 给出 CT18 的主要结果——作为 $x$ 和 $\Q$ 函数的拟合 PDF、QCD 参数（$\alpha_s$、$m_c$）的确定、计算的部分子光度，以及各种 PDF 矩和求和规则。这些比较将令广大将在 LHC 实验的理论预言中使用 PDF 的研究人员感兴趣。

Sec.~\ref{sec:Quality} 描述 CT18 对所拟合数据的成功理论描述能力。除刻画单个数据集的拟合质量外，Sec.~\ref{sec:StandardCandles} 还计算了与 LHC 唯象学相关的各类标准烛光量，例如 13 与 14 TeV 的希格斯玻色子产生截面，以及电弱玻色子与正反顶夸克对产生截面之间的各种关联。Sec.~\ref{sec:Conclusions} 讨论本工作更广泛的意义并突出我们的主要结论。

若干附录给出了重要的支撑细节。Appendix~\ref{sec:AppendixCT18Z} 综述 CT18Z 及其他替代拟合，包括进入这些独立分析的各数据集的描述。与似然函数及协方差矩阵之间关系有关的较形式化的细节概括于 Appendix~\ref{sec:chi2_app}。Appendix~\ref{sec:AppendixParam} 给出 CT18 采用的解析拟合形式及 PDF 参数的最优拟合值。Appendix~\ref{sec:AppendixCodeDevelopment} 给出 CT 拟合框架的多项技术进展，包括代码并行化；Appendix~\ref{sec:ATLASjetdecorrel} 列出拟合新纳入的 LHC 单举喷注数据时所用的去关联模型。最后，Appendix~\ref{sec:Appendix4xFitter} 给出基于 Hessian 剖析方法的简短研究的结果，用以评估 ATLAS 7 TeV $W/Z$ 产生数据的影响。
